\documentclass[11pt,a4paper]{article}
\usepackage[utf8]{inputenc}
\usepackage{kotex}
\usepackage[scale=0.80]{geometry}
\usepackage{hyperref}
\usepackage{parskip}

\begin{document}

\begin{flushright}
\textbf{Prabhavit Gupta}\\
Venice, Italy\\
+39 351 610 7807\\
prabhavitg@gmail.com\\
\href{https://www.linkedin.com/in/prabhavit/}{LinkedIn} | \href{https://github.com/prab-gupta}{GitHub}
\end{flushright}

\vspace{2em}

\textbf{To:} 채용 담당자님, 마이클 (Mycle)\\
\textbf{Date:} \today\\
\textbf{Subject:} 백엔드 엔지니어 (서비스 개발) 지원

\vspace{1em}

안녕하세요,

마이클(Mycle)의 백엔드 엔지니어 (서비스 개발) 포지션에 지원하게 되어 기쁩니다. 저는 복잡한 도메인의 문제를 해결하고, 높은 정합성과 안정성을 요구하는 시스템을 설계하는 데 강한 열정을 가진 백엔드 개발자입니다. 모빌리티 시장 혁신을 이끌어가는 마이클의 비전과, 개발팀이 지향하는 주도적이고 AI 친화적인 엔지니어링 문화에 깊이 공감하여 지원하게 되었습니다.

저는 FlowSpace에서 백엔드 엔지니어로 근무하며 이벤트 기반 아키텍처(Event-driven Architecture)와 비동기 작업 처리 시스템을 설계하고 확장한 경험이 있습니다. 복잡한 API 연동 과정에서 트랜잭션 경계를 설정하고, 웹훅(Webhook)의 재시도 로직 및 엣지 케이스를 처리하여 예약 도메인과 같이 정합성이 중요한 환경에서 시스템의 결함 허용성(Fault Tolerance)을 보장했습니다. 또한 데이터베이스의 쿼리를 튜닝하고 RDBMS 스키마를 최적화하여 서비스 성능을 대폭 개선한 경험이 있습니다.

제 주요 기술 스택은 Python, Node.js 기반이지만 새로운 환경에 매우 빠르게 적응합니다. 특히 실무에서 \textbf{Claude Code}와 같은 AI 코딩 도구를 적극적으로 활용하여 레거시 코드의 맥락을 파악하고 단순 반복 작업을 자동화하며 개발 생산성을 극대화하고 있습니다. 이러한 경험을 바탕으로 마이클의 Kotlin/Java 및 Spring 환경에도 신속하게 온보딩할 것이며, 팀 내에서 AI 도구를 일관되고 효율적으로 사용할 수 있도록 워크플로를 정비하는 데 기여하고 싶습니다.

저는 기능 개발 시 단순히 요청받은 것을 구현하는 것을 넘어, "이 기능이 왜 필요한가", "더 나은 대안은 없는가"를 치열하게 고민하며 동료들과 코드 리뷰를 통해 함께 성장하는 환경을 즐깁니다. 마이클팀에 합류하여 안정적이고 확장 가능한 백엔드 시스템을 구축하고 서비스 혁신에 기여하고 싶습니다.

제 지원서를 검토해 주셔서 감사합니다. 저는 현재 이탈리아에 거주 중이나 본 포지션을 위해 한국으로 인력 배치(Relocation) 및 원격 근무가 가능합니다. 면접을 통해 제 경험이 마이클의 서비스 발전에 어떻게 기여할 수 있을지 논의할 수 있기를 기대합니다.

감사합니다.

\vspace{1em}
Prabhavit Gupta

\end{document}
